\section{Data preparation \label{sec:data_prep}}


\textbf{Questionnaires}

To compare the introductory questionnaire and the final questionnaire, we created a composite index (design\_index) in the feedback questionnaire by calculating the average responses of the following questions:

\textbf{Playful variant:} (What is your attitude towards playful study materials?)
\begin{itemize}
    \item I would enjoy if other study materials were more playful as well
    \item The playful design (memes, badges, language) motivated me
    \item I enjoyed the playful narration in the theory section
\end{itemize}
\textbf{Competitive variant:} (What is your attitude towards competitive activities?)
\begin{itemize}
    \item I enjoyed the aggressive narration in the theory section
    \item I would enjoy if I were to compete with my classmates more often
    \item Seeing my performance compared to others motivated me
\end{itemize}

% \todo{tento satisfaction index... nebude len robit bordel v korelaciach? Nemali by sme ho len vypocitat raz radsej?}
We created another index (satisfaction\_index) that numerically expresses the usefulness of creating similar e-learning materials. It was calculated as the average response to the following questions:
\begin{itemize}
    \item I enjoyed the format of the course	
    \item I would voluntarily take a similar course.
    \item The course kept me engaged.
    \item This e-learning material added value to the university course PSI.
    \item I would like to encounter more e-learning materials during my studies at FIIT STU.
\end{itemize}

We substituted responses to the question \enquote{Was any of the topics of the course insufficiently explained or hard to understand?} with the number of topics the student marked.

\textbf{Reports and logs}

We combined the logs and reports from the previous section into a single table for better clarity and processing. This was only applied to reports that contained data for individual students - data about exercises, categories and questionnaire answers were not included. 
\todo{link to appendix}

% -Derived features:
%     -theory/quiz interactions ratio
%     -restructure interactions so that they are in buckets
%     -theory clicks

% -Data bucketing
%     -duration cat
%     -(interactions)

% -cleaning
%     -time was capped at 2:15 hrs per user
%     -removed feedback interactions
% -disregarded reports

    
%     # we have theory pages in C: 3+10+10+3=26
% # theory pages in P: 3+9+12+3=27

Time taken to complete quizzes in the course was capped at 135 minutes. The reason is that students could leave a quiz open and leave, resulting in course attempts spanning even days. We only included time spent on quiz attempts, as total time spent on the courses was too blurred by students' breaks to be interpreted.

The resulting dataset contained features describing each student's performance, behavior and derived metrics. The resulting features were:

\begin{itemize}
    \item id - user id for the participants' identification
    \item finalgrade - final grade calculated from weighted quiz results
    \item intro\_score - score in the introductory assessment quiz 
    \item final\_score - score in the final assessment quiz
    \item improvement - final\_score-intro\_score
    \item duration\_mins - time taken completing quizzes
    \item duration\_cat - time taken completing quizzes expressed divided into 4 categories 
    \item total\_attempts - total number of attempts at quizzes
    \item quiz -  number of interactions with quizzes
    \item theory - number of interactions with the theory pages
    \item theory/quiz interactions ratio
    \item extra\_theory - number of theory pages subtracted from all theory interactions
    \item extra\_theory/quiz interactions ratio
    \item extra\_attempts - attempts at quizzes beyond the mandatory one attempt per quiz
    \item extra\_theory\_per\_minute - extra\_theory/duration\_mins
    \item course - P or C, indicating which course the data point belongs to
\end{itemize}

\textbf{Exercises}

The data needed very little cleaning. The most difficult exercises made the categories they belonged to appear more difficult than they were. We attribute this to exercise design rather than insufficient understanding of the concepts. The exercises were removed from their respective categories, but kept for overall aggregate measurements. 

Additional dataset was created that calculates the average performance across categories per course type.